% ============================================================================
%  diagramas.tex — biblioteca TikZ do documento
%
%  Todos os diagramas do documento (C4, UML de casos de uso, atividades,
%  sequencia, entidade-relacionamento, classes e infraestrutura) sao vetoriais
%  e desenhados nativamente em TikZ. Isso garante tipografia unica (Roboto),
%  nitidez em qualquer zoom e ausencia de dependencia externa de renderizacao.
%  O codigo Mermaid equivalente de cada diagrama esta no Apendice A.
% ============================================================================

\usepackage{tikz}
\usetikzlibrary{
  arrows.meta, positioning, fit, backgrounds, calc, shapes.geometric,
  shapes.symbols, shapes.multipart, shapes.misc, decorations.pathreplacing,
  decorations.pathmorphing, decorations.markings, patterns, patterns.meta,
  matrix, chains, scopes, quotes, shadows.blur, math, intersections
}

\usepackage{pgfplots}
\pgfplotsset{compat=1.18}
\usepgfplotslibrary{groupplots}

% Numeracao dos graficos na convencao brasileira (virgula decimal), coerente
% com o restante do documento.
\pgfkeys{/pgf/number format/.cd, use comma, 1000 sep={.}}

% Estilo comum dos graficos quantitativos: eixos recessivos, grade sutil,
% tipografia do documento. Mesma paleta validada da interface.
\pgfplotsset{
  unooplot/.style={
    width=\linewidth, height=6cm,
    axis line style={draw=unooAxis, line width=0.6pt},
    tick label style={font=\scriptsize\color{unooInkSoft}},
    label style={font=\small\color{unooInkSoft}},
    title style={font=\small\bfseries\color{unooNavy}, align=left},
    legend style={font=\scriptsize, draw=unooGrid, fill=unooSurface,
                  /tikz/every even column/.append style={column sep=6pt}},
    grid=major, grid style={draw=unooGrid, line width=0.5pt},
    major tick style={draw=unooAxis},
    axis background/.style={fill=unooSurface},
    every axis plot post/.append style={line width=0.9pt},
    enlarge x limits=0.06,
  }
}

% ---------------------------------------------------------------------------
%  Base comum
% ---------------------------------------------------------------------------
\tikzset{
  >={Stealth[length=5.4pt, width=4.2pt]},
  every node/.append style={align=center, inner sep=4pt},
  cxlabel/.style={font=\scriptsize\color{unooInkSoft}, fill=unooSurface,
                  inner sep=1.6pt, rounded corners=1pt},
  cxtitle/.style={font=\small\bfseries\color{unooNavy}},
  legenda/.style={font=\scriptsize\color{unooMuted}},
}

% ---------------------------------------------------------------------------
%  C4 model — notacao de Simon Brown adaptada a paleta do projeto
% ---------------------------------------------------------------------------
\tikzset{
  % --- elementos -----------------------------------------------------------
  c4box/.style={
    rectangle, rounded corners=2.5pt, draw, line width=0.8pt,
    text width=32mm, minimum height=15mm, font=\footnotesize,
    inner sep=5pt
  },
  c4person/.style={
    c4box, fill=unooNavy, draw=unooNavy!85!black, text=white,
    rounded corners=7pt, text width=28mm, minimum height=16mm
  },
  c4system/.style={
    c4box, fill=unooBlue, draw=unooNavy, text=white
  },
  c4systemext/.style={
    c4box, fill=unooMuted, draw=unooInkSoft, text=white
  },
  c4container/.style={
    c4box, fill=unooBlue!88, draw=unooNavy, text=white
  },
  c4containerext/.style={
    c4box, fill=unooMuted!85, draw=unooInkSoft, text=white
  },
  c4component/.style={
    c4box, fill=unooBlueLight!45, draw=unooBlue, text=unooInk, text width=30mm,
    minimum height=13mm
  },
  c4db/.style={
    cylinder, shape border rotate=90, aspect=0.28, draw, line width=0.8pt,
    fill=unooNavy, draw=unooNavy!80!black, text=white,
    text width=25mm, minimum height=17mm, font=\footnotesize, inner sep=3pt
  },
  c4queue/.style={
    c4box, fill=unooYellow!85, draw=unooYellow!60!black, text=unooInk,
    rounded corners=8pt
  },
  c4artifact/.style={
    c4box, fill=unooGreen!80, draw=unooGreen!60!black, text=white
  },
  % --- fronteiras ----------------------------------------------------------
  c4boundary/.style={
    draw=unooAxis, dashed, dash pattern=on 3pt off 2.4pt, line width=0.9pt,
    rounded corners=4pt, inner sep=9pt
  },
  c4boundarylbl/.style={
    font=\scriptsize\bfseries\color{unooInkSoft}, anchor=north west,
    inner sep=2pt
  },
  c4zone/.style={
    c4boundary, draw=unooBlue!55, fill=unooBlue!4
  },
  c4zoneext/.style={
    c4boundary, draw=unooMuted, fill=unooPaper
  },
  % --- relacoes ------------------------------------------------------------
  c4rel/.style={-{Stealth[length=5pt,width=4pt]}, line width=0.85pt,
                draw=unooInkSoft},
  c4relbi/.style={c4rel, {Stealth[length=5pt,width=4pt]}-{Stealth[length=5pt,width=4pt]}},
  c4reldash/.style={c4rel, dashed, dash pattern=on 3pt off 2pt},
  c4relhi/.style={c4rel, draw=unooOrange, line width=1.1pt},
}

% Rotulo tecnologico dentro de um elemento C4
\newcommand{\tech}[1]{\\[1.5pt]{\scriptsize\itshape [#1]}}
% Descricao curta dentro de um elemento C4
\newcommand{\desc}[1]{\\[2pt]{\scriptsize #1}}

% ---------------------------------------------------------------------------
%  UML — casos de uso
% ---------------------------------------------------------------------------
\tikzset{
  ucase/.style={
    ellipse, draw=unooBlue, line width=0.85pt, fill=unooBlueWash,
    text=unooInk, font=\scriptsize, minimum width=34mm, minimum height=11mm,
    inner sep=2.5pt, text width=31mm
  },
  ucasesec/.style={ucase, fill=unooSurface, draw=unooAxis},
  ucsystem/.style={
    draw=unooNavy, line width=1pt, rounded corners=3pt, fill=unooBlue!3,
    inner sep=11pt
  },
  ucrel/.style={line width=0.8pt, draw=unooInkSoft},
  ucinclude/.style={
    -{Stealth[length=5pt,width=4pt]}, dashed, dash pattern=on 3pt off 2pt,
    line width=0.75pt, draw=unooGreen
  },
  ucextend/.style={
    -{Stealth[length=5pt,width=4pt]}, dashed, dash pattern=on 3pt off 2pt,
    line width=0.75pt, draw=unooOrange
  },
  ucgen/.style={-{Triangle[open, length=6pt, width=6pt]}, line width=0.8pt,
                draw=unooInkSoft},
  ucnote/.style={font=\scriptsize\color{unooInkSoft}},
}

% Ator UML (boneco de palito) — #1 nome do no, #2 posicao, #3 rotulo
\newcommand{\ator}[3]{%
  \begin{scope}[shift={#2}]
    \draw[line width=0.9pt, draw=unooNavy] (0,0.30) circle (0.155);
    \draw[line width=0.9pt, draw=unooNavy] (0,0.145) -- (0,-0.24);
    \draw[line width=0.9pt, draw=unooNavy] (-0.245,-0.01) -- (0.245,-0.01);
    \draw[line width=0.9pt, draw=unooNavy] (0,-0.24) -- (-0.20,-0.60);
    \draw[line width=0.9pt, draw=unooNavy] (0,-0.24) -- (0.20,-0.60);
    \node[font=\scriptsize\bfseries\color{unooNavy}, text width=24mm,
          anchor=north] at (0,-0.68) {#3};
    \node[inner sep=0pt, minimum width=13mm, minimum height=19mm] (#1) at (0,-0.15) {};
  \end{scope}
}

% Ator secundario (sistema externo) — retangulo com estereotipo
\newcommand{\atorsys}[3]{%
  \node[draw=unooInkSoft, line width=0.85pt, fill=unooPaper, rounded corners=2pt,
        font=\scriptsize, text width=22mm, minimum height=12mm] (#1) at #2
        {{\scriptsize$\ll$ator$\gg$}\\\textbf{#3}};
}

% ---------------------------------------------------------------------------
%  UML — diagramas de atividades
% ---------------------------------------------------------------------------
\tikzset{
  actstart/.style={
    circle, fill=unooInk, draw=unooInk, minimum size=4.6mm, inner sep=0pt
  },
  actend/.style={
    circle, fill=unooInk, draw=unooInk, minimum size=4.6mm, inner sep=0pt,
    double, double distance=1.1pt, outer sep=0pt
  },
  action/.style={
    rectangle, rounded corners=4pt, draw=unooBlue, line width=0.85pt,
    fill=unooBlueWash, text=unooInk, font=\scriptsize, text width=36mm,
    minimum height=10mm, inner sep=4pt
  },
  actionsys/.style={action, fill=unooSurface, draw=unooAxis},
  actionerr/.style={action, fill=unooOrangeWash, draw=unooOrange},
  actionok/.style={action, fill=unooGreenWash, draw=unooGreen},
  decision/.style={
    diamond, draw=unooNavy, line width=0.85pt, fill=unooYellowWash,
    aspect=1.9, font=\scriptsize, text width=26mm, inner sep=1pt
  },
  bar/.style={
    rectangle, fill=unooInk, minimum width=34mm, minimum height=1.5mm,
    inner sep=0pt
  },
  barv/.style={
    rectangle, fill=unooInk, minimum width=1.5mm, minimum height=30mm,
    inner sep=0pt
  },
  flow/.style={-{Stealth[length=5pt,width=4pt]}, line width=0.85pt,
               draw=unooInkSoft},
  flowlbl/.style={font=\scriptsize\color{unooInkSoft}, fill=unooSurface,
                  inner sep=1.5pt},
  lane/.style={
    draw=unooAxis, line width=0.7pt, fill=unooSurface
  },
  lanehdr/.style={
    draw=unooAxis, line width=0.7pt, fill=unooNavy, text=white,
    font=\scriptsize\bfseries
  },
  actnote/.style={
    draw=unooYellow!70!black, fill=unooYellowWash, line width=0.6pt,
    font=\scriptsize\color{unooInkSoft}, text width=34mm, align=left,
    rounded corners=0pt, inner sep=4pt
  },
  actnoteline/.style={dashed, dash pattern=on 2pt off 1.6pt, draw=unooYellow!70!black},
  objflow/.style={
    rectangle, draw=unooGreen, line width=0.7pt, fill=unooGreenWash,
    font=\scriptsize, text width=30mm, inner sep=3.5pt
  },
}

% ---------------------------------------------------------------------------
%  UML — diagramas de sequencia
% ---------------------------------------------------------------------------
\tikzset{
  seqhead/.style={
    rectangle, rounded corners=2pt, draw=unooNavy, line width=0.8pt,
    fill=unooNavy, text=white, font=\scriptsize, text width=26mm,
    minimum height=9mm, inner sep=3pt
  },
  seqheadext/.style={seqhead, fill=unooMuted, draw=unooInkSoft},
  seqheadact/.style={seqhead, fill=unooBlue, draw=unooNavy},
  seqline/.style={dashed, dash pattern=on 2.6pt off 2.2pt, draw=unooAxis,
                  line width=0.7pt},
  seqmsg/.style={-{Stealth[length=5pt,width=4pt]}, line width=0.85pt,
                 draw=unooInkSoft},
  seqret/.style={seqmsg, dashed, dash pattern=on 2.6pt off 2pt,
                 draw=unooInkSoft},
  seqself/.style={seqmsg, rounded corners=2pt},
  seqlbl/.style={font=\scriptsize\color{unooInk}, fill=unooSurface,
                 inner sep=1.6pt},
  seqact/.style={fill=unooBlueLight!60, draw=unooBlue, line width=0.5pt,
                 minimum width=2.2mm, inner sep=0pt},
  seqframe/.style={draw=unooAxis, line width=0.7pt, rounded corners=0pt},
  seqframelbl/.style={fill=unooNavy, text=white, font=\scriptsize\bfseries,
                      inner sep=2.2pt, anchor=north west},
  seqcond/.style={font=\scriptsize\itshape\color{unooInkSoft}, anchor=west},
}

% Raia de sequencia: #1 = x, #2 = profundidade (y final, negativo),
% #3 = titulo, #4 = tecnologia
\newcommand{\raia}[4]{%
  \node[seqhead] at (#1,0) {\textbf{#3}\\[0.5pt]{\scriptsize\itshape #4}};
  \draw[seqline] (#1,-0.52) -- (#1,#2);
}
\newcommand{\raiaext}[4]{%
  \node[seqheadext] at (#1,0) {\textbf{#3}\\[0.5pt]{\scriptsize\itshape #4}};
  \draw[seqline] (#1,-0.52) -- (#1,#2);
}
% Mensagem: #1 = x origem, #2 = x destino, #3 = y, #4 = rotulo
\newcommand{\msgs}[4]{%
  \draw[seqmsg] (#1,#3) -- (#2,#3) node[seqlbl, midway, above=0.4pt]{#4};
}
% Retorno (tracejado)
\newcommand{\msgr}[4]{%
  \draw[seqret] (#1,#3) -- (#2,#3) node[seqlbl, midway, above=0.4pt]{#4};
}
% Auto-mensagem: #1 = x, #2 = y, #3 = rotulo
\newcommand{\msgself}[3]{%
  \draw[seqself] (#1,#2) -- ++(0.62,0) -- ++(0,-0.42) -- ++(-0.62,0);
  \node[seqlbl, anchor=west] at (#1+0.72,#2-0.21) {#3};
}

% ---------------------------------------------------------------------------
%  Entidade-relacionamento e classes
% ---------------------------------------------------------------------------
\tikzset{
  entidade/.style={
    rectangle split, rectangle split parts=2, draw=unooNavy,
    line width=0.85pt, rectangle split part fill={unooNavy, unooSurface},
    font=\scriptsize, text width=37mm, inner sep=4pt, rounded corners=1pt,
    every text node part/.style={align=center},
    every two node part/.style={align=left}
  },
  classe/.style={
    rectangle split, rectangle split parts=3, draw=unooNavy,
    line width=0.85pt, rectangle split part fill={unooNavy, unooSurface, unooPaper},
    font=\scriptsize, text width=44mm, inner sep=4pt, rounded corners=1pt,
    every text node part/.style={align=center},
    every two node part/.style={align=left},
    every three node part/.style={align=left}
  },
  erel/.style={line width=0.85pt, draw=unooInkSoft},
  ercard/.style={font=\scriptsize\color{unooNavy}, fill=unooSurface,
                 inner sep=1.4pt},
  hered/.style={-{Triangle[open, length=6.5pt, width=6.5pt]}, line width=0.85pt,
                draw=unooInkSoft},
  compos/.style={-{Diamond[open, length=6pt, width=4.5pt]}, line width=0.85pt,
                 draw=unooInkSoft},
  depend/.style={-{Stealth[length=5pt,width=4pt]}, dashed,
                 dash pattern=on 3pt off 2pt, line width=0.75pt,
                 draw=unooInkSoft},
}

% Titulo branco em negrito para a primeira particao de entidade/classe
\newcommand{\etitulo}[1]{\textcolor{white}{\textbf{#1}}}
\newcommand{\pk}[1]{\textcolor{unooNavy}{\textbf{#1}}}
\newcommand{\fk}[1]{\textcolor{unooOrange}{#1}}

% ---------------------------------------------------------------------------
%  Infraestrutura / implantacao
% ---------------------------------------------------------------------------
\tikzset{
  nuvem/.style={
    draw=unooBlue, line width=1pt, dash pattern=on 4pt off 2.6pt,
    rounded corners=6pt, fill=unooBlue!3, inner sep=11pt
  },
  nuvemlbl/.style={
    font=\scriptsize\bfseries\color{unooNavy}, anchor=north west,
    fill=unooSurface, inner sep=2.4pt
  },
  regiao/.style={
    draw=unooAxis, line width=0.8pt, dash pattern=on 2.6pt off 2pt,
    rounded corners=4pt, fill=unooSurface, inner sep=8pt
  },
  servico/.style={
    rectangle, rounded corners=2.5pt, draw=unooNavy, line width=0.85pt,
    fill=unooBlue!88, text=white, font=\scriptsize, text width=30mm,
    minimum height=13mm, inner sep=4pt
  },
  servicoger/.style={servico, fill=unooNavy},
  servicobroker/.style={servico, fill=unooYellow!85, text=unooInk,
                        draw=unooYellow!55!black},
  volume/.style={
    cylinder, shape border rotate=90, aspect=0.30, draw=unooGreen!60!black,
    line width=0.85pt, fill=unooGreen!80, text=white, font=\scriptsize,
    text width=22mm, minimum height=15mm, inner sep=3pt
  },
  externo/.style={
    rectangle, rounded corners=2.5pt, draw=unooInkSoft, line width=0.85pt,
    fill=unooMuted, text=white, font=\scriptsize, text width=30mm,
    minimum height=13mm, inner sep=4pt
  },
  internet/.style={
    cloud, cloud puffs=13, cloud puff arc=110, draw=unooAxis, line width=0.85pt,
    fill=unooPaper, text=unooInkSoft, font=\scriptsize, minimum width=26mm,
    minimum height=15mm, aspect=1.6
  },
  netpriv/.style={-{Stealth[length=5pt,width=4pt]}, line width=0.9pt,
                  draw=unooGreen!70!black},
  netpub/.style={-{Stealth[length=5pt,width=4pt]}, line width=0.9pt,
                 draw=unooOrange},
  netbi/.style={{Stealth[length=5pt,width=4pt]}-{Stealth[length=5pt,width=4pt]},
                line width=0.9pt, draw=unooInkSoft},
  portlbl/.style={font=\scriptsize\ttfamily\color{unooInkSoft},
                  fill=unooSurface, inner sep=1.5pt},
}

% ---------------------------------------------------------------------------
%  Pipeline / fluxo de dados
% ---------------------------------------------------------------------------
\tikzset{
  etapa/.style={
    rectangle, rounded corners=2.5pt, draw=unooBlue, line width=0.85pt,
    fill=unooBlueWash, font=\scriptsize, text width=27mm, minimum height=13mm,
    inner sep=4pt
  },
  etapaio/.style={
    etapa, fill=unooPaper, draw=unooAxis
  },
  etapaout/.style={etapa, fill=unooGreenWash, draw=unooGreen},
  etapaerr/.style={etapa, fill=unooOrangeWash, draw=unooOrange},
  fita/.style={
    signal, signal to=east, signal from=west, draw=unooBlue,
    line width=0.85pt, fill=unooBlueWash, font=\scriptsize, text width=24mm,
    minimum height=12mm
  },
  pipe/.style={-{Stealth[length=5.4pt,width=4.4pt]}, line width=1pt,
               draw=unooBlue},
  pipelbl/.style={font=\scriptsize\color{unooInkSoft}, fill=unooSurface,
                  inner sep=1.5pt},
}

% ---------------------------------------------------------------------------
%  Legenda padronizada de diagrama
% ---------------------------------------------------------------------------
\newcommand{\amostra}[1]{%
  \tikz[baseline=-0.55ex]{\node[#1, text width=0pt, minimum width=3.6mm,
    minimum height=3.0mm, inner sep=0pt, font=\tiny]{};}%
}
\newenvironment{legendadiag}
  {\par\vspace{2pt}\begingroup\footnotesize\color{unooInkSoft}%
   \setlength{\parskip}{2pt}\textbf{Legenda:}\space}
  {\endgroup\par\vspace{2pt}}

% ---------------------------------------------------------------------------
%  Acrescimos deste documento (nao vieram da biblioteca de origem)
% ---------------------------------------------------------------------------
\tikzset{
  % Bundle estatico servido pela CDN — nao e um servico, e um artefato.
  bundle/.style={
    rectangle, draw=unooGreen!60!black, line width=0.85pt, fill=unooGreen!80,
    text=white, font=\scriptsize, text width=28mm, minimum height=13mm,
    inner sep=4pt, rounded corners=1pt
  },
  % Fronteira do iframe: a linha que separa o chrome do hub da ficcao da
  % experiencia. Tracejada e grossa de proposito — e fronteira de seguranca.
  iframe/.style={
    draw=unooOrange, line width=1.1pt, dash pattern=on 4pt off 2.4pt,
    rounded corners=3pt, fill=unooOrangeWash!45, inner sep=9pt
  },
  iframelbl/.style={
    font=\scriptsize\bfseries\color{unooOrange}, anchor=north west,
    fill=unooSurface, inner sep=2.2pt
  },
  % Maquina de estados
  estado/.style={
    rectangle, rounded corners=6pt, draw=unooBlue, line width=0.85pt,
    fill=unooBlueWash, font=\scriptsize, text width=24mm, minimum height=9mm,
    inner sep=3pt
  },
  estadoini/.style={estado, fill=unooNavy, text=white, draw=unooNavy},
  estadofim/.style={estado, fill=unooGreenWash, draw=unooGreen},
  estadoerr/.style={estado, fill=unooRedWash, draw=unooRed},
  trans/.style={-{Stealth[length=5pt,width=4pt]}, line width=0.8pt,
                draw=unooInkSoft},
  translbl/.style={font=\scriptsize\color{unooInkSoft}, fill=unooSurface,
                   inner sep=1.4pt},
  % Camada horizontal (diagramas de camadas)
  camada/.style={
    rectangle, draw=unooAxis, line width=0.7pt, fill=unooSurface,
    minimum height=13mm, inner sep=5pt, font=\scriptsize
  },
  camadalbl/.style={font=\scriptsize\bfseries\color{unooNavy}},
  % Tabela do banco, com prefixo por experiencia
  tabela/.style={
    rectangle split, rectangle split parts=2, draw=unooNavy,
    line width=0.8pt, rectangle split part fill={unooNavy, unooSurface},
    font=\scriptsize, text width=34mm, inner sep=3.5pt, rounded corners=1pt,
    every text node part/.style={align=center},
    every two node part/.style={align=left}
  },
  tabelaplat/.style={tabela, rectangle split part fill={unooBlue, unooSurface}},
  tabelacat/.style={tabela, rectangle split part fill={unooMuted, unooPaper}},
}

% Rotulo de prefixo de tabela por experiencia
\newcommand{\pref}[1]{{\scriptsize\ttfamily #1}}

% ---------------------------------------------------------------------------
%  UML 2.5.1 --- os tipos que a biblioteca de origem nao tinha (Secao
%  "Modelagem UML", secoes/05b-modelagem-uml.tex; ADR-122)
%
%  Pacote, componente com interfaces fornecidas e requeridas, no de
%  implantacao, artefato, objeto, estrutura composta (parte e porta), perfil,
%  fragmento de interacao, comunicacao e tempo. A notacao segue a
%  especificacao da OMG (UML 2.5.1, formal/17-12-05); a cor segue os tokens
%  unoo* como o resto do documento.
% ---------------------------------------------------------------------------
\newcommand{\estereo}[1]{{\guillemotleft}#1{\guillemotright}}

\tikzset{
  % --- relacoes de classe que faltavam ------------------------------------
  associa/.style={line width=0.85pt, draw=unooInkSoft},
  navega/.style={associa, -{Stealth[length=5pt,width=4pt]}},
  agrega/.style={associa, {Diamond[open, length=7.5pt, width=5pt]}-},
  compoe/.style={associa, {Diamond[length=7.5pt, width=5pt, fill=unooInkSoft]}-},
  realiza/.style={-{Triangle[open, length=6.5pt, width=6.5pt]}, dashed,
                  dash pattern=on 3pt off 2pt, line width=0.8pt, draw=unooInkSoft},
  % extensao de perfil: seta cheia, que nao se confunde com a heranca (aberta)
  estende/.style={-{Triangle[length=6.5pt, width=6.5pt, fill=unooInkSoft]},
                  line width=0.85pt, draw=unooInkSoft},
  multip/.style={font=\scriptsize\color{unooNavy}, inner sep=1.2pt},
  papelassoc/.style={font=\scriptsize\itshape\color{unooInkSoft}, inner sep=1.2pt},
  % --- classificadores ----------------------------------------------------
  interface/.style={classe, draw=unooLavenderSoft,
    rectangle split part fill={unooLavenderSoft, unooSurface, unooPaper}},
  enumeracao/.style={entidade, draw=unooInkSoft, text width=30mm,
    rectangle split part fill={unooMuted, unooSurface}},
  metaclasse/.style={rectangle, draw=unooInkSoft, line width=0.8pt,
    fill=unooPaper, font=\scriptsize, text width=28mm, minimum height=11mm,
    inner sep=3pt},
  % --- objeto: nome sublinhado, instancia e nao tipo -----------------------
  objeto/.style={rectangle split, rectangle split parts=2, draw=unooBlue,
    line width=0.85pt, rectangle split part fill={unooBlueLight!55, unooSurface},
    font=\scriptsize, text width=40mm, inner sep=3.5pt,
    every text node part/.style={align=center},
    every two node part/.style={align=left}},
  objetoum/.style={rectangle, draw=unooBlue, line width=0.85pt,
    fill=unooBlueLight!45, font=\scriptsize, text width=30mm,
    minimum height=10mm, inner sep=3pt},
  objetolink/.style={line width=0.8pt, draw=unooInkSoft},
  % --- pacote --------------------------------------------------------------
  pacotecorpo/.style={draw=unooNavy, line width=0.85pt, fill=unooBlueWash,
    inner sep=0pt},
  pacoteaba/.style={draw=unooNavy, line width=0.85pt, fill=unooNavy,
    text=white, font=\scriptsize\bfseries, inner sep=2.4pt,
    minimum height=4.6mm},
  pacotedep/.style={-{Stealth[length=5pt,width=4pt]}, dashed,
    dash pattern=on 3pt off 2pt, line width=0.8pt, draw=unooInkSoft},
  % --- componente: retangulo com o icone de componente no canto -------------
  componente/.style={rectangle, draw=unooNavy, line width=0.85pt,
    fill=unooBlueWash, font=\scriptsize, text width=30mm, minimum height=15mm,
    inner sep=4pt,
    path picture={
      \coordinate (icc) at ([xshift=-3.0mm, yshift=-3.1mm]path picture bounding box.north east);
      \draw[draw=unooNavy, line width=0.45pt, fill=unooSurface]
        ([xshift=-1.3mm, yshift=-1.8mm]icc) rectangle ([xshift=1.6mm, yshift=1.8mm]icc);
      \draw[draw=unooNavy, line width=0.45pt, fill=unooSurface]
        ([xshift=-2.1mm, yshift=0.5mm]icc) rectangle ([xshift=-0.5mm, yshift=1.2mm]icc);
      \draw[draw=unooNavy, line width=0.45pt, fill=unooSurface]
        ([xshift=-2.1mm, yshift=-1.2mm]icc) rectangle ([xshift=-0.5mm, yshift=-0.5mm]icc);
    }},
  componenteext/.style={componente, fill=unooPaper, draw=unooInkSoft},
  ifacelbl/.style={font=\scriptsize\color{unooInkSoft}, inner sep=1.4pt},
  % --- implantacao ---------------------------------------------------------
  nodocorpo/.style={draw=unooNavy, line width=0.85pt, fill=unooSurface,
    inner sep=0pt},
  nodolbl/.style={font=\scriptsize\color{unooNavy}, align=left, inner sep=3pt},
  artefato/.style={rectangle, draw=unooGreen!60!black, line width=0.8pt,
    fill=unooGreenWash, font=\scriptsize, text width=30mm, minimum height=11mm,
    inner sep=3.5pt,
    path picture={
      \coordinate (ica) at ([xshift=-2.9mm, yshift=-2.9mm]path picture bounding box.north east);
      \draw[draw=unooGreen!60!black, line width=0.45pt, fill=unooSurface]
        ([xshift=-1.3mm, yshift=1.7mm]ica) -- ([xshift=-1.3mm, yshift=-1.7mm]ica)
        -- ([xshift=1.3mm, yshift=-1.7mm]ica) -- ([xshift=1.3mm, yshift=0.8mm]ica)
        -- ([xshift=0.4mm, yshift=1.7mm]ica) -- cycle;
      \draw[draw=unooGreen!60!black, line width=0.45pt]
        ([xshift=0.4mm, yshift=1.7mm]ica) -- ([xshift=0.4mm, yshift=0.8mm]ica)
        -- ([xshift=1.3mm, yshift=0.8mm]ica);
    }},
  implanta/.style={-{Stealth[length=5pt,width=4pt]}, dashed,
    dash pattern=on 3pt off 2pt, line width=0.75pt, draw=unooGreen!60!black},
  caminhocom/.style={line width=0.9pt, draw=unooInkSoft},
  % --- estrutura composta: parte, porta e conector -------------------------
  parte/.style={rectangle, draw=unooNavy, line width=0.85pt, fill=unooSurface,
    font=\scriptsize, text width=28mm, minimum height=11mm, inner sep=3pt},
  partemult/.style={parte, fill=unooBlueWash},
  porta/.style={rectangle, draw=unooNavy, line width=0.8pt, fill=unooSurface,
    minimum size=2.6mm, inner sep=0pt},
  portaborda/.style={porta, fill=unooOrange!85, draw=unooOrange!70!black},
  conector/.style={line width=0.85pt, draw=unooInkSoft},
  % --- interacao: operador de fragmento e uso de interacao (ref) -----------
  frameop/.style={chamfered rectangle, chamfered rectangle corners=south east,
    chamfered rectangle xsep=1.4mm, chamfered rectangle ysep=1.2mm,
    fill=unooNavy, draw=unooNavy, text=white, font=\scriptsize\bfseries,
    inner sep=1.8pt, anchor=north west},
  guarda/.style={font=\scriptsize\itshape\color{unooInkSoft}, anchor=west,
    inner sep=1.4pt},
  refint/.style={rectangle, draw=unooNavy, line width=0.85pt, fill=unooSurface,
    font=\scriptsize, text width=40mm, minimum height=11mm, inner sep=3pt},
  % --- comunicacao ---------------------------------------------------------
  msgcomlbl/.style={font=\scriptsize\color{unooInk}, fill=unooSurface,
    inner sep=1.4pt},
  % --- maquina de estados: pseudoestados e estado composto -----------------
  escolha/.style={diamond, draw=unooNavy, line width=0.85pt, fill=unooSurface,
    minimum size=6mm, inner sep=0pt},
  estadocomposto/.style={draw=unooBlue, line width=0.9pt, rounded corners=8pt,
    fill=unooBlue!4},
  % --- tempo ---------------------------------------------------------------
  tempofaixa/.style={draw=unooAxis, line width=0.7pt, fill=unooSurface},
  tempolinha/.style={line width=1.2pt, draw=unooBlue!70!unooNavy},
  tempoestado/.style={font=\scriptsize\color{unooInkSoft}, anchor=east,
    inner sep=1.5pt},
  tempotick/.style={font=\scriptsize\color{unooMuted}, inner sep=1.5pt},
  tempomarca/.style={dashed, dash pattern=on 2pt off 1.6pt, draw=unooAxis,
    line width=0.6pt},
  restricao/.style={font=\scriptsize\color{unooNavy}, fill=unooSurface,
    inner sep=1.4pt},
}

% Titulo de classificador com estereotipo: #1 estereotipo, #2 nome (branco,
% para a primeira particao escura de classe/entidade/interface)
\newcommand{\ctitulo}[2]{\textcolor{white}{{\scriptsize\estereo{#1}}\\[-0.5pt]\textbf{#2}}}

% Pacote: #1 opcoes do corpo, #2 nome do no, #3 e #4 cantos opostos (fit),
% #5 rotulo da aba
\newcommand{\pacote}[5][]{%
  \node[pacotecorpo, fit={#3 #4}, #1] (#2) {};
  \node[pacoteaba, anchor=south west] at (#2.north west) {#5};
}

% No de implantacao (caixa em perspectiva): #1 opcoes do corpo, #2 nome do
% no, #3 e #4 cantos opostos (fit), #5 rotulo (estereotipo + nome)
\newcommand{\nodo}[5][]{%
  \node[nodocorpo, fit={#3 #4}, #1] (#2) {};
  \filldraw[draw=unooNavy, line width=0.85pt, fill=unooNavy!22]
    (#2.north west) -- ++(2.4mm,2.4mm) -- ([shift={(2.4mm,2.4mm)}]#2.north east)
    -- (#2.north east) -- cycle;
  \filldraw[draw=unooNavy, line width=0.85pt, fill=unooNavy!40]
    (#2.north east) -- ++(2.4mm,2.4mm) -- ([shift={(2.4mm,2.4mm)}]#2.south east)
    -- (#2.south east) -- cycle;
  \node[nodolbl, anchor=north west] at (#2.north west) {#5};
}

% Interface fornecida (pirulito): #1 centro da bola
\newcommand{\bola}[1]{%
  \draw[line width=0.85pt, draw=unooInkSoft, fill=unooSurface] #1 circle (1.3mm);}
% Interface requerida (soquete): #1 centro, #2 angulo para onde o soquete ABRE
\newcommand{\soquete}[2]{%
  \draw[line width=0.85pt, draw=unooInkSoft]
    #1 ++(#2+90:2.2mm) arc[start angle=#2+90, end angle=#2+270, radius=2.2mm];}

% Mensagem de diagrama de comunicacao: seta curta ao lado do link.
% #1 posicao do rotulo, #2 ponto medio, #3 angulo da seta, #4 rotulo
\newcommand{\msgcom}[4][above]{%
  \draw[seqmsg] #2 ++(#3+180:4.5mm) -- ++(#3:9mm)
    node[msgcomlbl, midway, #1=1.5pt]{#4};}
